\scanpage{399}\mspage{313}

est cofinale dans $I$. Comme c'est une sous-catégorie pleine (et $=$ str. pleine), il suffit de prouver la propriété F1, i.e. que pour tout $i\in\mathrm{Ob}\,I$, existe $i'\in\mathrm{Ob}\,I'$ \deleted{tel que} et une flèche $i\longrightarrow i'$.

On va construire un $i'$ par construction transfinie, sur l'intervalle ordinal $I_{\alpha_\pi}$, en construisant un système inductif

\[
(i_\alpha)_{\alpha\in I_{\alpha_\pi}}\qquad i_\alpha\in\mathrm{Ob}\,I
\]

d'objets de $I$, et \deleted{une famille} \inserted{un système} \deleted{pour tout $\alpha\in I_{\alpha_\pi}$ un $j_\alpha\in J'$ et une flèche $\unsure{\varphi|\phi}(i_\alpha)\longrightarrow j_\alpha$.} \inserted{inductif}

\[
\inserted{(j_\alpha)_{\alpha\in I_{\alpha_\pi}},\qquad j_\alpha\in\mathrm{Ob}\,J',}
\]

\deleted{\unsure{(systif)|(respectif)}} flèches de transition notées

\[
\tau_{\beta,\alpha}:i_\alpha\unsure{-|\longrightarrow}i_\beta,\qquad
\tau'_{\beta,\alpha}:j_\alpha\unsure{-|\longrightarrow}j_\beta,
\]

enfin pour tout $\alpha\unsure{\text{dans}|\in}I_{\unsure{<|\alpha}\pi}$ des flèches

\diagram{\begin{tikzcd}[ampersand replacement=\&,column sep=small]
(*)\quad \& \varphi(i_\alpha)\arrow[r,"u_\alpha"]\arrow[rr,bend right=24,dashed,"{\varphi(\tau_{\alpha+1,\alpha})}"']\arrow[rrr,no head,bend right=33,start anchor=south,end anchor=south] \& j_\alpha\arrow[r,"v_\alpha"]\arrow[rr,bend left=28,dashed,"{\tau'_{\alpha+1,\alpha}}"] \& \varphi(i_{\alpha+1})\,[\arrow[r,dotted,"u_{\alpha+1}"] \& j_{\alpha+1}\,]
\end{tikzcd}}

dans $J$, telles que

\[
\left\{\begin{aligned}
v_\alpha u_\alpha&=\varphi(\tau_{\alpha+1,\alpha})\unsure{=|,}\\
u_{\alpha+1}v_\alpha&=\tau'_{\alpha+1,\alpha},
\end{aligned}\right.\tag*{$(**)$}
\]

% slanted
\margin{J'aurais \unsure{(aussi)} pu \unsure{prendre}}

% archive sheet number: 398
